Das Prinzip der kapazitiven Füllstandsmessung beruht darauf, dass die Kapazität eines Kondensators vom Dielektrikum zwischen den Elektroden abhängt. In diesem Versuch wird diese Eigenschaft genutzt, um zu untersuchen, wie die Eintauchtiefe des Kondensators in ein Dielektrikum seine Kapazität verändert und wie aus der Kapazitätsänderung auf die Füllstandshöhe geschlossen werden kann.
In dieser Arbeit werden das Messprinzip eines kapazitiven Füllstandssensors, der Aufbau der Versuchsanordnung sowie der Ablauf der digitalen Signalverarbeitung analysiert. Dabei werden der Aufbau der Oszillatorschaltung, die digitale Erfassung der Messsignale sowie die Datenverarbeitung und -auswertung mit Matlab behandelt. Außerdem werden mögliche Fehlerquellen untersucht und die Messgenauigkeit bewertet.
